\section{Numerical method}\label{sec:method}

\status{Ready as to what is implemented. The shared-machinery argument is the
same as in the three-dimensional companion and is verified in the code, not
assumed. The two-dimensional specifics below --- the normal and curvature
estimator, and the form of the front --- were read from the source, and where
the two dimensions diverge that is stated rather than smoothed over.}

\subsection{Shared machinery, and the single point of difference}

The comparison in this paper is only meaningful if the two methods differ in
\emph{one} thing. They are therefore built on the same components: the same
hierarchical grid, the same interpolated properties, the same momentum solver,
the same projection, and coupling at the same point of the multiphase step.
What differs is how the interface is represented, and everything the
representation entails --- how the normal and the curvature are obtained, and
how the interface is transported.

\missing{the itemised correspondence between the two paths through the
multiphase step, as in the 3D companion, with the line numbers of the single
branch point.}

\subsection{VOF: geometric PLIC in the plane}

The volume fraction is advected by a directionally split geometric scheme: one
sweep per direction, each reconstructing a straight segment per interface cell
from the fraction in its neighbourhood, fluxing the area that segment cuts, and
copying the auxiliary field back before the next sweep.

The normal and the curvature come from a \emph{height-function} estimator. The
two-dimensional path uses the variant of Shirani; the three-dimensional
companion uses the standard height function with Mehta's correction. This is a
genuine difference between the two papers and not a presentational one: the
estimator is part of what is being compared, and quoting a 2D number against a
3D one without saying so would be comparing two different methods.

\caveat{The code also carries ELVIRA and an adaptive height-function variant,
both reachable and both currently disabled at the call site. The measurements
in this paper are with Shirani only.}

\subsection{Front-tracking: explicit polyline}

In the plane the front is a closed, ordered polyline of markers with explicit
connectivity, advected by the interpolated velocity and maintained by local
surgery --- splitting an edge that grows beyond a tolerance, merging one that
falls below it. Surface tension enters as a force distributed from the markers
to the grid.

\caveat{The two-dimensional front-tracking core carries a header comment
describing it as phase B1 --- kinematic advection in a prescribed field,
without the solver. The coupling is supplied by a separate adapter in
\texttt{examples-common}, and the example's build file describes the assembled
pair as phase~B2. Before any production claim is made from this path, the
degree of coupling actually exercised has to be stated precisely; the 301-frame
Newtonian run on disk indicates a coupled run, but that is an inference from
the output, not a reading of the call graph.}

\subsection{Interface extraction for comparison}

Both methods are measured through the \emph{same} quantities, computed by the
same code path: the integrals that define circularity, centre of mass and rise
velocity are taken over the volume-fraction field, which VOF advects directly
and which front-tracking writes from the geometry of the front. The measurement
is therefore identical on both sides and any difference that survives belongs
to the method.

\subsection{Adaptive refinement}

\missing{the two-dimensional criterion, as implemented: seeds from the
interface, dilation to the refinement levels, and the minimum-cell rule.}

The one result already established here is negative and worth stating early,
because it governs how any convergence study in this paper must be set up. With
the refined band specified as a fixed \emph{number of cells} per side, each
resolution refines a \emph{different physical width} --- the finest mesh
refines the largest fraction of the domain --- and a mesh-convergence study
built that way does not converge for reasons that have nothing to do with the
discretisation.

\missing{the measured evidence for the paragraph above: five series, none
monotone, with non-monotonicity between 12 and 40 times the sampling floor.
Port the numbers from the report and put them in \texttt{dados/}.}

The remedy is to anchor the band to a fixed \emph{physical} width, taken from
the coarsest resolution of the series, so that finer meshes only gain cells.
A floor on the number of fine cells per side is kept and \emph{overrides} the
physical band when the two disagree, because the reconstruction and the
fraction transport both need that neighbourhood: it is a correctness
constraint, not a tuning parameter.
